%% paper/sections/intro.tex
%% Draft only. No results, no methods, no conclusions -- see
%% docs/DECISION_RULES.md and logs/GATES.md for why those wait.

A text token has a referent that survives its position: the word ``teacher''
at sequence index 3 and the same word at index 40 name the same lexical
item, and a language transformer's attention over those two positions is
answering a well-posed question --- \emph{how related are these two
mentions of the same thing?} A patch token in a video transformer has no
such guarantee. A grid cell at spatial location $(i,j)$ in frame $t$ and the
same cell in frame $t{+}1$ denote the same thing only if nothing has moved
through that cell, an assumption video violates constantly and by design ---
people walk, cameras pan, occlusion happens. Divided space-time and
factorised attention over a fixed patch grid nonetheless price every pairwise
affinity between these positions as if the grid supplied a stable identity
axis. It does not. The tokens are referent-free.

This paper asks a narrow, checkable question that follows from that
observation: does giving a video transformer tokens with a genuine
persistent referent --- one token per tracked entity per time segment,
rather than one token per spatial cell --- change what the model can learn,
and is the effect attributable to \emph{persistence itself} rather than to
two things that usually travel with it, entity-localized cropping (looking
only at where an entity is, instead of the whole frame) and a smaller,
matched token budget? Prior entity- and object-centric video representations
motivate the persistence hypothesis but rarely isolate it from these
confounds (\S\ref{sec:related-work}); we design a control that shuffles
entity identity across time segments while holding architecture, token
count, and every other factor fixed, so that any gap between the shuffled
and unshuffled conditions is attributable to persistence and only to
persistence (\S\ref{sec:method}).

Answering that question honestly turned out to require answering a
prior one first: is the benchmark we asked it on actually testing what it
claims to test? Our first candidate primary dataset, OUC-CGE, a recent
classroom group-engagement corpus~\cite{lu2025ouccge}, failed that check. A
linear probe on a single mean-pooled frame embedding --- no temporal
structure, no entities, no attention, nothing an architecture paper would
want credit for --- reaches
%% outputs/phase3_5_audit/REPORT.md sec.4
macro-F1 $0.9748$ on its provided validation split, statistically
indistinguishable from a full entity-persistent transformer trained on the
same split (
%% outputs/phase3_5_audit/REPORT.md sec.4, citing the locked-recipe A1 result
$0.9767$). The same probe, unchanged, scores
%% outputs/dataset_admission/daisee/admission_report.json: trivial_probe.val_macro_f1
$0.2177$ on DAiSEE~\cite{gupta2016daisee}, a dataset built with a
subject-disjoint split. The gap is not a property of the probe; it is a
property of the split. \S\ref{sec:dataset-audit} diagnoses why, reframes the
finding as a property of how OUC-CGE's release was partitioned rather than a
criticism of its collection or its authors, and turns the diagnosis into a
reusable admission procedure that any candidate dataset can be checked
against before it is trusted with a hypothesis test. We apply that procedure
to our eventual primary dataset before using it, not after being burned by
it.

The contributions of this paper are therefore: (1) Entity-Persistent
Tokenization (EPT) and an identity-shuffle control that isolates persistence
from its usual confounds, evaluated under a decision rule and effect-size
floor pre-registered before any test-set evaluation
(\S\ref{sec:pre-registration}); (2) a diagnosed, quantified account of
train/test leakage in a released classroom-engagement benchmark, reframed
in \S\ref{sec:dataset-audit} as independent confirmation of a known failure
mode in clip-derived video benchmarks rather than a claim of discovering the
phenomenon itself; and (3) a lightweight, reusable dataset-admission test ---
cross-split near-duplicate detection plus a trivial-feature-probe floor
against a subject-disjoint reference --- that we run on every dataset in
this paper, not only the one that failed it.
